\chapter{Glossary}
\label{ch:glossary}

\begin{longtable}{@{}p{0.24\textwidth}p{0.70\textwidth}@{}}
\caption{Course glossary}\label{tab:glossary}\\
\toprule
Term & Meaning in this course \\
\midrule
\endfirsthead
\toprule
Term & Meaning in this course \\
\midrule
\endhead
\bottomrule
\endfoot
Analytical grain &
What one row represents in a relation or result. Also called the unit of
observation. \\
Archive &
One packaged file that contains other files, commonly in ZIP format. \\
Attempt &
One student's participation in one module presentation. A student can have
multiple attempts. \\
Audit &
A recorded inspection of data or files against stated expectations. \\
Candidate key &
A column or combination of columns expected to identify one row uniquely. It
must be tested. \\
Cardinality &
The number and pattern of relationships between rows, such as one-to-many. \\
Causal claim &
A claim about what would change because of an intervention or exposure, not
merely what is associated in observed data. \\
Code and script &
Code is text containing computer instructions. A script is a file that runs a
sequence of those instructions. \\
Command and terminal &
A command is one instruction given to a computer. A terminal is the
text-based application in which project commands are run. \\
Constraint &
A database rule that restricts permitted values or relationships. \\
Core layer &
The project's typed, named, analysis-ready relational layer. \\
CSV &
Comma-separated values: a text format for tabular data. \\
Data and dataset &
Data are recorded values or facts. A dataset is an organized collection of
records. \\
Data science &
The combined use of programming, databases, statistics, and subject knowledge
to learn from data and communicate evidence. \\
Data dictionary &
Documentation describing fields, values, tables, and their intended meanings. \\
Database &
An organized collection of data managed by a database system. In PostgreSQL, a
database contains schemas, which contain relations and other objects. \\
Descriptive analysis &
Analysis that summarizes what is recorded in a defined dataset and population. \\
Dimension &
A descriptive attribute used to group or filter measures, such as module or
presentation. \\
Evidence log &
A record connecting a result to its query, unit, population, interpretation,
and limitations. \\
Feature &
A variable constructed for analysis or modeling. Its definition must include
grain and information-availability timing. \\
Foreign key &
A value or set of values that refers to a key in another relation. \\
Grain &
See analytical grain. \\
Identifier &
A value used to name or match a recorded entity. Whether it uniquely identifies
a row depends on the table's grain. \\
IMD band &
A binned index-of-multiple-deprivation field in OULAD. It contains missing
values and should not be treated as a complete measure of an individual. \\
Information leakage &
Use of data that would not be available at the intended decision time, or that
otherwise reveals the target improperly. \\
Implementation &
The working code and settings that perform a project, as distinct from
documentation that explains it. \\
Input and output &
An input is what a step receives; an output is what the step produces. \\
Join &
An operation that combines rows from relations according to a match condition. \\
Lineage &
The traceable path from original source data through transformations to a
reported result. \\
Materialized view &
A stored result of a query that must be refreshed when its inputs change. \\
Metric &
A precisely defined value used to summarize something, including its unit and
denominator where applicable. \\
Missing value &
An absent or unavailable field value. Its meaning depends on source and
context; it is not automatically an error or zero. \\
Module &
An anonymized OULAD course unit. \\
Model &
A mathematical rule used to summarize relationships or calculate predictions
from input variables. \\
Normalization &
Organizing relations to reduce redundancy and update anomalies. \\
Null &
SQL's marker for a missing or unknown value. Comparisons use
\texttt{IS NULL}, not equality to null. \\
Observation window &
The time interval from which information is permitted when constructing an
analysis or prediction. \\
OULAD &
The Open University Learning Analytics Dataset used by this project. \\
Orphan &
A child row whose expected parent key cannot be found. \\
Path &
The address of a file or directory, such as
\projectfile{sql/02_staging.sql}. \\
Pipeline &
An ordered workflow in which the output of one step becomes input to another. \\
Population &
The complete set of observations about which an analytical claim is made. \\
PostgreSQL &
The open-source relational database system used by the project. \\
Presentation &
A particular offering of a module, identified by an anonymized code. \\
Primary key &
The database-enforced identifier for a row in a relation. \\
Prediction &
An estimate of a later or unknown recorded outcome. It does not by itself
explain why the outcome occurs. \\
Proportion &
A part divided by a named whole. Multiplying a proportion by 100 produces a
percentage. \\
Proxy &
A recorded stand-in for a concept that is difficult to measure directly. \\
Quality gate &
A set of blocking checks that must pass before dependent analysis proceeds. \\
Query &
A request written in SQL that asks a database to return, combine, summarize,
or modify data. \\
Raw layer &
The source-preserving database layer, loaded primarily as text so source
conditions remain inspectable. \\
Reconciliation &
Comparison of totals or other invariants across source, raw, transformed, and
reported stages. \\
Referential integrity &
Consistency of relationships between child and parent keys. \\
Relation &
A table-like set of rows and columns in the relational model. \\
Repository &
The organized project folder and its tracked files. \\
Reproducibility &
The ability to regenerate results from documented source, code,
configuration, and environment. \\
Row and column &
A row records one observation at the table's grain. A column stores one kind
of value across rows. \\
Schema &
A namespace inside a PostgreSQL database. This project uses schemas to separate
raw, staging, core, and quality responsibilities. \\
Staging layer &
The conversion layer where raw text is parsed, typed, normalized, and checked
before publication to core relations. \\
Surrogate key &
An identifier created by a system rather than derived from the natural domain
attributes. \\
Unit of analysis &
The entity about which a result is calculated or interpreted. It must agree
with the query grain. \\
Unlogged table &
A PostgreSQL table that avoids write-ahead logging for speed but has weaker
crash-recovery guarantees. \\
Validation &
Testing data against stated expectations while preserving evidence of any
failure. \\
Variable &
A recorded characteristic that can differ across observations and is usually
stored in a table column. \\
VLE &
Virtual learning environment. OULAD records summarized interaction with VLE
sites by student and relative day. \\
Warning profile &
A nonblocking quality result that requires an analytical decision or
interpretive note. \\
Window function &
A SQL function that calculates across related rows without collapsing them in
the way \texttt{GROUP BY} does. \\
\end{longtable}
